\section{Preliminaries}\label{sec:prelim}

Throughout this paper, $E$ denotes a Banach space over $\R$ with norm
$\|\cdot\|_E$. In all concrete examples, $E=\R^d$ for some positive integer
$d$. For a bounded linear operator $T$ on~$E$, we write $\|T\|$ for the
operator norm, $\sigma(T)$ for the spectrum, and $\rho(T)=\max\{|\lambda|:
\lambda\in\sigma(T)\}$ for the spectral radius. We denote by $I$ the identity
operator on~$E$ and by $0$ the zero operator.

\subsection{Reasoning trajectory spaces}

We model a reasoning process as a finite sequence of states in~$E$, where each
state represents a $d$-dimensional embedding of a reasoning step.

\begin{definition}[Trajectory space]\label{def:trajectory}
  Let $n\geq 1$ be a positive integer. The \emph{trajectory space of length~$n$}
  is $\mathcal{M}_n = E^{n+1}$, equipped with the normalized $\ell^1$ metric
  \begin{equation}\label{eq:metric}
    d(\tau,\tau') = \frac{1}{n+1}\sum_{i=0}^{n}\|s_i - s_i'\|_E,
  \end{equation}
  where $\tau=(s_0,s_1,\dots,s_n)$ and $\tau'=(s_0',s_1',\dots,s_n')$.
  An element $\tau\in\mathcal{M}_n$ is called a \emph{reasoning trajectory};
  $s_0$ is the \emph{initial query} and $(s_1,\dots,s_n)$ are the
  \emph{reasoning steps}.
\end{definition}

Since $E$ is complete, the product space $(\mathcal{M}_n,d)$ is a complete
metric space.

\subsection{CoT operators}

\begin{definition}[Linear CoT operator]\label{def:cot}
  A \emph{linear CoT operator} $T=T(A,B)$, parameterized by matrices
  $A,B\in\R^{d\times d}$, acts on trajectories by
  \begin{equation}\label{eq:cot-action}
    T(A,B)(s_0,s_1,\dots,s_n) = (s_0,\, As_0+Bs_1,\, As_1+Bs_2,\,
    \dots,\, As_{n-1}+Bs_n).
  \end{equation}
  The matrix $A$ is the \emph{context propagation matrix} and $B$ is the
  \emph{current-state transformation matrix}.
\end{definition}

The operator $T(A,B)$ preserves the initial query~$s_0$ and transforms each
subsequent state $s_i$ ($i\geq 1$) using a linear combination of the previous
state $s_{i-1}$ (weighted by~$A$) and the current state $s_i$ (weighted
by~$B$). This captures the essential structure of CoT reasoning: each step
combines contextual information from the preceding step with a refinement of
the current reasoning state.

\begin{definition}[Block matrix representation]\label{def:block}
  The CoT operator $T(A,B)$ on $\mathcal{M}_n$ admits the block matrix
  representation $M(A,B)\in\R^{(n+1)d\times(n+1)d}$:
  \begin{equation}\label{eq:block-matrix}
    M(A,B) =
    \begin{pmatrix}
      I & 0 & 0 & \cdots & 0 \\
      A & B & 0 & \cdots & 0 \\
      0 & A & B & \cdots & 0 \\
      \vdots & & \ddots & \ddots & \vdots \\
      0 & 0 & \cdots & A & B
    \end{pmatrix},
  \end{equation}
  where each entry is a $d\times d$ block.
\end{definition}

\begin{definition}[Sub-operator spectral radius]\label{def:rho-sub}
  For a CoT operator $T(A,B)$, the \emph{sub-operator spectral radius} is
  $\rho_{\mathrm{sub}}(T) = \rho(B)$, the spectral radius of the current-state
  transformation matrix~$B$.
\end{definition}

\subsection{Semigroup theory prerequisites}

We recall several classical results that we use in our proofs. For background
on operator semigroups, we refer to~\cite{Gluck2018} and~\cite{Gerlach2017}.

\begin{definition}[Semigroup and monoid]\label{def:semigroup}
  A \emph{semigroup} is a non-empty set $S$ equipped with an associative binary
  operation. A \emph{monoid} is a semigroup with an identity element.
\end{definition}

\begin{definition}[Semigroup at infinity]\label{def:semigroup-infty}
  For a bounded operator semigroup $\mathcal{T}=\{T_s\}_{s\in S}$ on a Banach
  space~$E$, the \emph{semigroup at infinity} is
  \[
    \mathcal{T}_\infty = \bigcap_{s\in S}
    \overline{\{T_t : t\geq s\}}^{\,\SOT},
  \]
  the set of all strong operator topology cluster points of the
  net~$(T_s)$~\cite{Gluck2018}.
\end{definition}

We use the following classical results without proof.

\begin{theorem}[Banach fixed-point theorem]\label{thm:banach}
  Let $(X,d)$ be a complete metric space and let $f\colon X\to X$ be a
  contraction with Lipschitz constant $\alpha<1$. Then $f$ has a unique fixed
  point $x^*\in X$, and for every $x_0\in X$,
  $d(f^n(x_0),x^*)\leq\frac{\alpha^n}{1-\alpha}\,d(x_0,f(x_0))$.
\end{theorem}

\begin{theorem}[JdLG decomposition {\cite{Gluck2018}}]\label{thm:jdlg-prereq}
  Let $\mathcal{T}$ be a bounded operator semigroup on a Banach space~$E$ such
  that $\mathcal{T}_\infty$ is non-empty and strongly compact. Then there
  exists a unique minimal idempotent $P_\infty\in\mathcal{T}_\infty$, and $E$
  decomposes as
  \[
    E = E_{\mathrm{rev}}\oplus E_{\mathrm{s}},
  \]
  where $E_{\mathrm{rev}}=\im(P_\infty)$ is the \emph{reversible part} on which
  $\mathcal{T}_\infty$ acts as a compact topological group, and
  $E_{\mathrm{s}}=\ker(P_\infty)$ is the \emph{stable part} on which every
  element of $\mathcal{T}_\infty$ acts as zero.
\end{theorem}

\begin{theorem}[Mean ergodic theorem]\label{thm:mean-ergodic}
  Let $T$ be a power-bounded operator on a reflexive Banach space~$E$. Then the
  Ces\`aro averages $\frac{1}{N}\sum_{k=0}^{N-1}T^k$ converge in the strong
  operator topology to the projection onto $\ker(I-T)$ along
  $\overline{\im(I-T)}$.
\end{theorem}

\begin{theorem}[Gelfand spectral radius formula]\label{thm:gelfand}
  For any bounded linear operator $T$ on a Banach space,
  $\rho(T)=\lim_{n\to\infty}\|T^n\|^{1/n}$.
\end{theorem}
